\documentclass[11pt]{article}
\usepackage[margin=1in]{geometry}
\usepackage{amsmath,amssymb}
\usepackage[colorlinks=true,linkcolor=blue,urlcolor=blue,citecolor=blue]{hyperref}
\setlength{\parindent}{0pt}
\setlength{\parskip}{0.7em}

\title{The rational-spectrum bottleneck in dimension one\\
\large An exact later-literature answer}
\author{Run \texttt{fa\_banach\_001} --- \texttt{agent\_lane\_02}}
\date{11 August 2026}

\begin{document}
\maketitle

\textbf{Status.} \texttt{literature\_already\_answered}. This is a
literature-identification note, not a claim of a new proof.

\section*{Source question}

Dutkay and Lai, \emph{Some reductions of the spectral set conjecture to
integers}, arXiv:1301.0814v2, page 4, Question (Q1), ask:
\begin{quote}
Is it true that every spectral set of Lebesgue measure $1$ has some rational
spectrum?
\end{quote}
The same issue is introduced on page 3, immediately after Remark 1.4, as the
missing property needed to reverse the implications among the spectral-to-tile
statements on $\mathbb R$, $\mathbb Z$, and all finite cyclic groups
$\mathbb Z_n$.

\section*{Exact supporting answer}

Fu and Song, \emph{Rational Spectra for Finite Hadamard Pairs and Bounded
Spectral Sets on the real line}, arXiv:2607.19858v1, prove on page 4:
\begin{quote}
\textbf{Theorem 1.3.} Let $\Omega\subset\mathbb R$ be a bounded measurable
spectral set with $|\Omega|=1$, and let $\Lambda$ be a spectrum of $\Omega$
with $0\in\Lambda$. Then $\Lambda\subset\mathbb Q$.
\end{quote}
The proof is on page 14. Their Remark 5.1 observes that $0\in\Lambda$ is only a
normalization: for any $\lambda_0\in\Lambda$, the translate
$\Lambda-\lambda_0$ is again a spectrum. Thus every spectral set in the source
question has a rational normalized spectrum. This is stronger than the
requested existence statement.

The supporting authors explicitly know they are answering the source question.
Their introduction cites Dutkay--Lai, describes rationality as the missing
hypothesis in that reduction, and states immediately after Theorem 1.3 that the
theorem supplies it. The mechanism is periodicity and fiberization to a finite
spectral pair, followed by their finite rationality theorem (Theorem 1.2).

\section*{Search, prior gap, and scope}

The local run indexes had no prior entry for arXiv:1301.0814 or the rational
spectrum question. Exact-title, phrase, and arXiv searches located the 2019
preprint arXiv:1908.02794 and the 2026 answer. The 2019 preprint was withdrawn
after a gap was found in Section 4, so it is not used as the supporting answer.
Fu--Song explicitly discuss that failed claim and replace its critical step
with a Galois/modulus-rigidity argument.

This identification closes Question (Q1) only. The full one-dimensional
Fuglede conjecture remains open, reduced to its finite cyclic analogues, and
Question (Q2) concerning the Coven--Meyerowitz property is not answered here.

\section*{References}

\begin{enumerate}
\item D. E. Dutkay and C.-K. Lai, \emph{Some reductions of the spectral set
conjecture to integers}, arXiv:1301.0814v2 (2013).
\item X.-Y. Fu and Z.-J. Song, \emph{Rational Spectra for Finite Hadamard
Pairs and Bounded Spectral Sets on the real line}, arXiv:2607.19858v1 (2026).
\item C.-K. Lai and Y. Wang, \emph{Spectrum is rational in dimension one},
arXiv:1908.02794v2 (withdrawn, 2020).
\end{enumerate}

\end{document}

